\documentclass[12pt,a4paper]{article}

% ── Packages ──────────────────────────────────────────────────────────────────
\usepackage[a4paper, margin=2.5cm]{geometry}
\usepackage{graphicx}
\usepackage{booktabs}
\usepackage{array}
\usepackage{longtable}
\usepackage{hyperref}
\usepackage{xcolor}
\usepackage{titlesec}
\usepackage{fancyhdr}
\usepackage{caption}
\usepackage{subcaption}
\usepackage{amsmath}
\usepackage{enumitem}
\usepackage{parskip}
\usepackage{float}
\usepackage{multirow}
\usepackage[T1]{fontenc}
\usepackage{lmodern}

% ── Colours ───────────────────────────────────────────────────────────────────
\definecolor{primary}{RGB}{31,78,121}
\definecolor{accent}{RGB}{68,114,196}
\definecolor{lightgray}{RGB}{245,245,245}

% ── Hyperref setup ────────────────────────────────────────────────────────────
\hypersetup{
    colorlinks=true,
    linkcolor=primary,
    urlcolor=accent,
    citecolor=primary,
    pdftitle={Driver Warning System — Project Report},
    pdfauthor={Aashish Biswocarma}
}

% ── Header / Footer ───────────────────────────────────────────────────────────
\pagestyle{fancy}
\fancyhf{}
\rhead{\textcolor{primary}{\small Driver Warning System}}
\lhead{\textcolor{primary}{\small FDS Project Report}}
\cfoot{\thepage}
\renewcommand{\headrulewidth}{0.4pt}

% ── Section title styling ─────────────────────────────────────────────────────
\titleformat{\section}{\large\bfseries\color{primary}}{}{0em}{}[\titlerule]
\titleformat{\subsection}{\normalsize\bfseries\color{accent}}{}{0em}{}

% ── Document ──────────────────────────────────────────────────────────────────
\begin{document}

% ── Title Page ────────────────────────────────────────────────────────────────
\begin{titlepage}
    \centering
    \vspace*{1.5cm}

    {\Huge\bfseries\color{primary} Driver Warning System}\\[0.4cm]
    {\Large\color{accent} Using Nepal Road Accident Data}\\[0.3cm]
    \rule{\linewidth}{1.5pt}\\[0.3cm]
    {\large\itshape Fundamentals of Data Science — Project Report}\\[2cm]

    \begin{tabular}{rl}
        \textbf{Author:}      & Aashish Biswocarma \\[0.3cm]
        \textbf{Repository:}  & \href{https://github.com/aashu2-prog/FDS_Project-IOE}{github.com/aashu2-prog/FDS\_Project-IOE} \\[0.3cm]
        \textbf{Date:}        & March 2026 \\[0.3cm]
    \end{tabular}

    \vfill
    \rule{\linewidth}{0.5pt}\\[0.2cm]
    {\small Institute of Engineering (IOE)}
\end{titlepage}

% ── Table of Contents ─────────────────────────────────────────────────────────
\tableofcontents
\newpage

% ══════════════════════════════════════════════════════════════════════════════
\section{Business Understanding}
% ══════════════════════════════════════════════════════════════════════════════

\subsection{Problem Statement}
Road accidents remain one of the leading causes of fatalities and injuries in Nepal. Drivers often lack real-time awareness of the risk level associated with their current location and driving conditions. This project aims to build a \textbf{Driver Warning System} that analyses a driver's location, road environment, and contextual factors to predict the likely severity of an accident and issue an appropriate safety warning.

\subsection{Objectives}
\begin{itemize}[leftmargin=1.5cm]
    \item Analyse the Nepal Road Accident dataset using the full Data Science lifecycle.
    \item Clean and preprocess the raw, deteriorated dataset to a model-ready state.
    \item Identify key factors that influence accident severity through Exploratory Data Analysis.
    \item Train a machine learning model to classify accident severity into four categories.
    \item Deploy a prototype warning system that maps a prediction to an actionable driver alert.
\end{itemize}

\subsection{Target Variable}
\texttt{Accident\_Severity} — four classes:

\begin{center}
\begin{tabular}{clc}
\toprule
\textbf{Label} & \textbf{Class} & \textbf{Warning Level} \\
\midrule
0 & Damage Only    & {\color{green!60!black} No warning} \\
1 & Minor Injury   & {\color{yellow!60!black} Caution} \\
2 & Serious Injury & {\color{orange} Drive carefully} \\
3 & Fatal          & {\color{red} High-risk zone} \\
\bottomrule
\end{tabular}
\end{center}

% ══════════════════════════════════════════════════════════════════════════════
\section{Data Acquisition}
% ══════════════════════════════════════════════════════════════════════════════

\subsection{Dataset Overview}
The raw dataset (\texttt{data/01\_raw/dataset.csv}) contains historical road accident records from Nepal spanning 2023–2026. It was intentionally deteriorated to simulate real-world data quality issues for demonstration purposes.

\begin{center}
\begin{tabular}{ll}
\toprule
\textbf{Property} & \textbf{Value} \\
\midrule
File            & \texttt{dataset.csv} \\
Raw rows        & 1,240 \\
Columns         & 15 \\
Target column   & \texttt{Accident\_Severity} \\
\bottomrule
\end{tabular}
\end{center}

\subsection{Features}
\begin{longtable}{lll}
\toprule
\textbf{Column} & \textbf{Type} & \textbf{Description} \\
\midrule
\endhead
Accident\_ID       & Categorical & Unique record identifier \\
Accident\_Date     & Date        & Date of accident \\
Province           & Categorical & Province (7 provinces) \\
District           & Categorical & District within province \\
Road\_Type         & Categorical & Highway / Urban / Rural / Feeder \\
Intersection\_Type & Categorical & Straight / Curve / Junction / Roundabout \\
Time\_of\_Day      & Categorical & Morning / Afternoon / Evening / Night \\
Weather\_Condition & Categorical & Clear / Rain / Fog / Snow \\
Road\_Surface      & Categorical & Dry / Wet / Slippery / Icy \\
Vehicle\_Type      & Categorical & Car / Bus / Truck / Motorbike etc. \\
Driver\_Gender     & Categorical & Male / Female / Other \\
Speed\_Zone        & Numeric     & Speed limit (km/h) \\
Year               & Numeric     & Year of accident \\
Accident\_Severity & Categorical & \textbf{Target} \\
\bottomrule
\end{longtable}

\subsection{Detected Quality Issues}
The following issues were identified during data acquisition:

\begin{enumerate}
    \item \textbf{Missing values} — Weather\_Condition (151), Road\_Surface (136), Driver\_Gender (114), District (90), Time\_of\_Day (73), Accident\_Date (10)
    \item \textbf{Duplicate rows} — 20+ exact duplicate records
    \item \textbf{Label typos} — e.g.\ \textit{Fatall}, \textit{minor injury}, \textit{Serious injury}, \textit{damage only}
    \item \textbf{Inconsistent casing} — Province values in ALL CAPS (e.g.\ \textit{BAGMATI})
    \item \textbf{Outliers} — Speed\_Zone values of 999 km/h
    \item \textbf{Year anomaly} — Year = 1999 (out of expected 2023–2026 range)
    \item \textbf{Corrupted dates} — Accident\_Date = `\texttt{N/A}' (non-parseable strings)
\end{enumerate}

% ══════════════════════════════════════════════════════════════════════════════
\section{Data Cleaning}
% ══════════════════════════════════════════════════════════════════════════════

Each quality issue was addressed systematically in notebook \texttt{03\_data\_cleaning.ipynb}.

\begin{center}
\begin{tabular}{lll}
\toprule
\textbf{Issue} & \textbf{Fix Applied} & \textbf{Result} \\
\midrule
Missing categoricals & Mode imputation per column    & 0 missing values \\
Duplicate rows       & \texttt{drop\_duplicates()}   & 33 rows removed \\
Severity typos       & String normalisation + mapping & 4 clean classes \\
Province casing      & \texttt{str.title()} transform & Consistent casing \\
Speed\_Zone outliers & Replace 999 with column median & Valid range only \\
Year anomaly (1999)  & Replace with mode year         & 2023–2026 only \\
Corrupted dates      & Parse errors → forward fill    & All dates valid \\
\bottomrule
\end{tabular}
\end{center}

\textbf{Before cleaning:} 1,240 rows \quad\textbf{After cleaning:} 1,207 rows

Output saved to \texttt{data/03\_interim/accidents\_cleaned.csv}.

% ══════════════════════════════════════════════════════════════════════════════
\section{Exploratory Data Analysis}
% ══════════════════════════════════════════════════════════════════════════════

\subsection{Target Distribution}
The dataset is moderately imbalanced. Minor Injury accounts for nearly half of all records.

\begin{center}
\begin{tabular}{lrr}
\toprule
\textbf{Class} & \textbf{Count} & \textbf{\%} \\
\midrule
Minor Injury   & 551 & 45.9\% \\
Serious Injury & 289 & 24.1\% \\
Damage Only    & 256 & 21.3\% \\
Fatal          & 104 &  8.7\% \\
\midrule
\textbf{Total} & \textbf{1,200} & \textbf{100\%} \\
\bottomrule
\end{tabular}
\end{center}

\subsection{Key Findings}

\begin{itemize}[leftmargin=1.5cm]
    \item \textbf{Province risk:} Koshi Province has the highest high-severity rate (36.1\%), followed by Bagmati (34.3\%).
    \item \textbf{High-risk districts:} Lamjung (53.6\%), Jhapa (45.5\%), and Chitwan (43.5\%) have the highest proportions of fatal/serious accidents.
    \item \textbf{Time of day:} Evening and Night show higher proportions of serious outcomes.
    \item \textbf{Road type:} Highway roads have the highest fatal rate (9.5\%) compared to Feeder roads (6.3\%).
    \item \textbf{Speed zone:} Serious Injury accidents cluster at higher speed zones (median 80 km/h vs 50 km/h for minor).
    \item \textbf{Day of week:} Friday records the most accidents; Severity is roughly uniform across weekdays.
    \item \textbf{Driver gender:} 76.5\% of accident-involved drivers are male across all provinces.
    \item \textbf{Vehicle type:} Jeep and Car show slightly higher fatal proportions than Microbus/Scooter.
\end{itemize}

\subsection{Correlation with Target}
Label-encoded features show low but non-zero correlation with severity, which is expected for nominal categorical variables. \texttt{District\_enc} and \texttt{Province\_enc} carry the highest predictive signal.

% ══════════════════════════════════════════════════════════════════════════════
\section{Feature Engineering}
% ══════════════════════════════════════════════════════════════════════════════

All categorical columns were transformed to numeric using \texttt{sklearn}'s \texttt{LabelEncoder}. Non-predictive columns (\texttt{Accident\_ID}, \texttt{Accident\_Date}) were dropped.

\begin{center}
\begin{tabular}{ll}
\toprule
\textbf{Column} & \textbf{Encoding} \\
\midrule
Province           & LabelEncoder \\
District           & LabelEncoder \\
Road\_Type         & LabelEncoder \\
Intersection\_Type & LabelEncoder \\
Time\_of\_Day      & LabelEncoder \\
Weather\_Condition & LabelEncoder \\
Road\_Surface      & LabelEncoder \\
Vehicle\_Type      & LabelEncoder \\
Driver\_Gender     & LabelEncoder \\
Speed\_Zone        & Numeric (kept as-is) \\
Year               & Numeric (kept as-is) \\
\midrule
\textbf{Target}    & \texttt{Severity\_Label} $\in \{0,1,2,3\}$ \\
\bottomrule
\end{tabular}
\end{center}

Output: \texttt{data/04\_processed/accidents\_features.csv} — 1,207 rows $\times$ 12 columns.

\textbf{Class distribution (full dataset):}
0: 258 \quad 1: 554 \quad 2: 291 \quad 3: 104

% ══════════════════════════════════════════════════════════════════════════════
\section{Modelling}
% ══════════════════════════════════════════════════════════════════════════════

\subsection{Model Choice}
A \textbf{Random Forest Classifier} was chosen for its robustness to noise, built-in feature importance, and interpretability — appropriate for a student project.

\begin{center}
\begin{tabular}{ll}
\toprule
\textbf{Hyperparameter} & \textbf{Value} \\
\midrule
\texttt{n\_estimators}  & 100 \\
\texttt{class\_weight}  & \texttt{`balanced'} \\
\texttt{random\_state}  & 42 \\
\bottomrule
\end{tabular}
\end{center}

\subsection{Train / Test Split}
80\% train / 20\% test, stratified by class.

\begin{center}
\begin{tabular}{lrr}
\toprule
\textbf{Split} & \textbf{Samples} & \textbf{\%} \\
\midrule
Train & 965 & 80\% \\
Test  & 242 & 20\% \\
\bottomrule
\end{tabular}
\end{center}

\subsection{Driver Warning Demo}
After prediction, the system maps the class label to a human-readable alert:

\begin{center}
\begin{tabular}{cll}
\toprule
\textbf{Predicted Class} & \textbf{Warning} & \textbf{Action} \\
\midrule
0 — Damage Only    & No warning          & Continue normally \\
1 — Minor Injury   & Caution advised     & Drive attentively \\
2 — Serious Injury & Drive carefully     & Reduce speed \\
3 — Fatal          & High-risk zone      & Avoid if possible \\
\bottomrule
\end{tabular}
\end{center}

% ══════════════════════════════════════════════════════════════════════════════
\section{Evaluation}
% ══════════════════════════════════════════════════════════════════════════════

\subsection{Overall Accuracy}
\[
\text{Test Accuracy} = 44.21\%
\]

The moderate accuracy is expected given: (a) small dataset size, (b) high class overlap in label-encoded categorical features, and (c) inherent difficulty of predicting accident severity from location/condition data alone.

\subsection{Classification Report}

\begin{center}
\begin{tabular}{lcccc}
\toprule
\textbf{Class} & \textbf{Precision} & \textbf{Recall} & \textbf{F1-Score} & \textbf{Support} \\
\midrule
Damage Only    & 0.25 & 0.13 & 0.17 & 52  \\
Minor Injury   & 0.49 & 0.82 & 0.61 & 111 \\
Serious Injury & 0.33 & 0.16 & 0.21 & 58  \\
Fatal          & 0.00 & 0.00 & 0.00 & 21  \\
\midrule
Weighted Avg   & 0.36 & 0.44 & 0.37 & 242 \\
\bottomrule
\end{tabular}
\end{center}

\textbf{Observations:}
\begin{itemize}[leftmargin=1.5cm]
    \item The model performs best on \textit{Minor Injury} (F1 = 0.61), the majority class.
    \item \textit{Fatal} class is never predicted — the dataset has only 104 fatal samples, making it hard to learn.
    \item Precision drops for minority classes due to class imbalance.
\end{itemize}

\subsection{ROC-AUC Scores (One-vs-Rest)}

\begin{center}
\begin{tabular}{lc}
\toprule
\textbf{Class} & \textbf{AUC} \\
\midrule
Damage Only    & 0.48 \\
Minor Injury   & 0.55 \\
Serious Injury & 0.62 \\
Fatal          & 0.54 \\
\bottomrule
\end{tabular}
\end{center}

\textit{Serious Injury} achieves the best AUC (0.62), suggesting the model partially distinguishes it from other classes via location and road type features.

\subsection{Feature Importances}

\begin{center}
\begin{tabular}{lc}
\toprule
\textbf{Feature} & \textbf{Importance} \\
\midrule
District\_enc           & 0.170 \\
Vehicle\_Type\_enc      & 0.115 \\
Province\_enc           & 0.114 \\
Road\_Type\_enc         & 0.098 \\
Time\_of\_Day\_enc      & 0.094 \\
Year                    & 0.085 \\
Road\_Surface\_enc      & 0.083 \\
Intersection\_Type\_enc & 0.075 \\
Speed\_Zone             & 0.074 \\
Weather\_Condition\_enc & 0.062 \\
\bottomrule
\end{tabular}
\end{center}

Geographic features (\textit{District}, \textit{Province}) dominate — confirming that location is the strongest predictor of accident severity in this dataset.

% ══════════════════════════════════════════════════════════════════════════════
\section{Conclusion}
% ══════════════════════════════════════════════════════════════════════════════

\subsection{Summary}
This project demonstrated the complete Data Science lifecycle — from raw, dirty data to a deployed prototype warning system. The key takeaways are:

\begin{itemize}[leftmargin=1.5cm]
    \item Data cleaning reduced the dataset from 1,240 to 1,207 rows and resolved 7 quality issues.
    \item EDA revealed that location (district/province), road type, and time of day are the most informative factors.
    \item A Random Forest model achieves 44.21\% accuracy — reasonable for a 4-class problem with limited, label-encoded features.
    \item The warning demo successfully maps model predictions to actionable driver alerts.
\end{itemize}

\subsection{Limitations}
\begin{itemize}[leftmargin=1.5cm]
    \item Small dataset (1,207 records) limits model generalisation.
    \item Label encoding of high-cardinality categoricals (e.g.\ District) loses ordinal meaning.
    \item Fatal class is heavily underrepresented (8.7\%), causing the model to ignore it.
\end{itemize}

\subsection{Future Work}
\begin{itemize}[leftmargin=1.5cm]
    \item Apply SMOTE oversampling to address class imbalance.
    \item Use frequency/target encoding instead of label encoding for high-cardinality features.
    \item Integrate a real-time GPS feed and serve predictions via a REST API.
    \item Experiment with Gradient Boosting (XGBoost/LightGBM) for better accuracy.
    \item Collect more data, especially for the Fatal class.
\end{itemize}

% ══════════════════════════════════════════════════════════════════════════════
\section*{Repository}
% ══════════════════════════════════════════════════════════════════════════════
\addcontentsline{toc}{section}{Repository}

All notebooks, data, models, and figures are available at:\\[0.2cm]
\href{https://github.com/aashu2-prog/FDS_Project-IOE}{\texttt{https://github.com/aashu2-prog/FDS\_Project-IOE}}

\begin{center}
\begin{tabular}{ll}
\toprule
\textbf{Path} & \textbf{Contents} \\
\midrule
\texttt{notebooks/01\_business\_understanding/} & Problem framing \\
\texttt{notebooks/02\_data\_acquisition/}        & Load \& inspect raw data \\
\texttt{notebooks/03\_data\_cleaning/}           & Fix all 7 quality issues \\
\texttt{notebooks/04\_eda/}                      & 22 visualisations \\
\texttt{notebooks/05\_feature\_engineering/}     & Encoding \& feature selection \\
\texttt{notebooks/06\_modeling/}                 & Train Random Forest + demo \\
\texttt{notebooks/07\_evaluation/}               & Metrics, ROC, feature importance \\
\texttt{data/01\_raw/}                           & Raw deteriorated dataset \\
\texttt{data/03\_interim/}                       & Cleaned dataset \\
\texttt{data/04\_processed/}                     & Encoded features \\
\texttt{models/trained/}                         & Serialised \texttt{.pkl} model \\
\texttt{reports/figures/}                        & All chart PNGs \\
\bottomrule
\end{tabular}
\end{center}

\end{document}
